% ===================================================================
%  図表第 6 号 — 家の中、家具、食事、明り。
%
%  Ceci n'est PAS une transcription : c'est une traduction, faite en
%  2026, en japonais d'aujourd'hui. Voir texto/ja/10-zuhyo-01.tex
%  pour la consigne, qui vaut pour les seize fichiers.
% ===================================================================

%%K t06-apar-1 apar
\VUsaut{6.0mm}
\VUtitre{15.0pt}{60}{図表第 6 号}
\VUsaut{1.4mm}
\VUtitre{11.4pt}{0}{家の中、家具、食事、}
\VUsaut{0.4mm}
\VUtitre{11.4pt}{0}{明り。}
\VUsaut{2.2mm}

%%K t06-apar-2 apar
家はできあがった。ジャンの叔父は新しい住まいに移った。その間取りは我々もう
知っている。では、どう飾りつけたかを見よう。まず見に行くのは次の部屋である。

%%K t06-tit-1 sub
\VUpk{10.2pt}{40}{寝室。}
\VUsaut{3.0mm}

%%K t06-c1-01-1 p
1. --- 我々は折わるく来てしまった。女中が\VUgras{部屋}を片づけている最中で
ある。

%%K t06-c1-02-1 p
2. --- \VUgras{洗面台} \textsuperscript{(1)} の上には、体を洗い、髪を梳き、
つまりは身支度をするための品があちこちに置いてある。\VUgras{洗面器}
\textsuperscript{(2)} と\VUgras{水差} \textsuperscript{(3)}、\VUgras{髪刷子}
\textsuperscript{(4)}、\VUgras{櫛} \textsuperscript{(5)}、\VUgras{石鹸}
\textsuperscript{(6)}、\VUgras{海綿} \textsuperscript{(7)}、それに最後に口を漱ぐ薬の
小さな\VUgras{瓶} \textsuperscript{(8)} である。

%%K t06-c1-03-1 p
3. --- 床の上、洗面台の前には汚れた水を受ける\VUgras{手桶} \textsuperscript{(9)}
がある。

%%K t06-c1-04-1 p
4. --- \VUgras{次の間} \textsuperscript{(10)} には\VUgras{衣鉤} \textsuperscript{(13)}
が幾つかと、\VUgras{傘} \textsuperscript{(11)} が二本見える。傘は\VUgras{傘立}
\textsuperscript{(12)} に挿してある。

%%K t06-c1-05-1 p
5. --- 戸の向う側には\VUgras{鏡付の箪笥} \textsuperscript{(14)} が立っていて、
中に\VUgras{衣類} \textsuperscript{(15)} が入っている。

%%K t06-c1-06-1 p
6. --- \VUgras{女中} \textsuperscript{(16)} が\VUgras{寝台} \textsuperscript{(17)} を
整えているうちに、その部分部分を見てゆこう。\VUgras{寝台の枠}
\textsuperscript{(20)} の上によい\VUgras{ばね蒲団} \textsuperscript{(21)} があり、
ジュスチーヌはやがてその上に毛の\VUgras{蒲団} \textsuperscript{(22)} を敷くので
ある。それから椅子の上の二枚の\VUgras{敷布} \textsuperscript{(23)} と
\VUgras{毛布} \textsuperscript{(24)} を取る。次に\VUgras{長枕} \textsuperscript{(25)}
と\VUgras{枕} \textsuperscript{(26)} を寝台の頭に置く。それらはいま\VUgras{羽蒲団}
\textsuperscript{(27)} とともに\VUgras{寝台前の敷物} \textsuperscript{(32)} の上に
ある。\VUgras{寝台の帳} \textsuperscript{(19)} と\VUgras{天蓋} \textsuperscript{(18)}
を羽根の払子で払えば、仕事は半ばできる。

%%K t06-c1-07-1 p
7. --- それから\VUgras{寝台脇の台} \textsuperscript{(28)} を少し片づけ、
\VUgras{目覚時計} \textsuperscript{(29)} に鍵を巻き、\VUgras{夜の灯}
\textsuperscript{(30)} を用意し、\VUgras{蝋燭} \textsuperscript{(31)} を持って行って
燭台を磨かねばならない。

%%K t06-tit-2 sub
\VUsaut{5.0mm}
\VUpk{10.2pt}{40}{針箱と暖炉の道具。}
\VUsaut{3.0mm}

%%K t06-c1-08-1 p
8. --- \VUgras{針箱} \textsuperscript{(34)} は\VUgras{床几} \textsuperscript{(33)} の
そばに置いてあり、これもたいそう片づけを要する。\VUgras{刺繍}
\textsuperscript{(35)} をたたみ、\VUgras{針仕事の台} \textsuperscript{(38)} に
\VUgras{指抜}、\VUgras{糸} \textsuperscript{(36)}、\VUgras{針} \textsuperscript{(37)}、
\VUgras{鋏} \textsuperscript{(39)} をしまい、台の蓋を\VUgras{鍵} \textsuperscript{(40)}
で閉めねばならない。

%%K t06-c1-09-1 p
9. --- 部屋の真中の\VUgras{丸卓} \textsuperscript{(41)} では、ジュスチーヌが
\VUgras{硝子瓶} \textsuperscript{(42)} の水を替える。\VUgras{砂糖壺}
\textsuperscript{(43)} に\VUgras{砂糖}を入れ、\VUgras{砂糖挟} \textsuperscript{(44)}
を磨き、\VUgras{衣刷子} \textsuperscript{(46)} を持って行く。

%%K t06-c1-10-1 p
10. --- 部屋の右手は仕事が少ない。\VUgras{長椅子} \textsuperscript{(47)} と
その\VUgras{座蒲団} \textsuperscript{(48)} はもとの場所にあり、寒くもないので
\VUgras{暖炉} \textsuperscript{(49)} の道具は誰も動かしていない。\VUgras{火掻}
\textsuperscript{(50)}、\VUgras{火挟} \textsuperscript{(51)}、\VUgras{薪架}
\textsuperscript{(55)}、\VUgras{灰よけ} \textsuperscript{(56)} はきれいである。
\VUgras{火よけ屏} \textsuperscript{(54)} も動いていないし、\VUgras{薪箱}
\textsuperscript{(52)} には\VUgras{割木} \textsuperscript{(53)} がいっぱいである。

%%K t06-noto-1 noto
\VUnotes{143.2mm}{%
(*) Babo = 小さな子供。英語 \textit{baby}、フランス語 \textit{bébé}、
イタリア語 \textit{bambino}。%
}

%%K t06-noto-2 noto
\VUnotes{143.2mm}{%
(*) Lambrequino = フランス語 \textit{lambrequin}、スペイン語
\textit{lambrequines}、ポルトガル語 \textit{lambrequins}、ドイツ語と英語まで
\textit{lambrequins} と言う --- つまり独、英、仏、葡、西の五つの言語で同じ
である。%
}

%%K t06-c1-11-1 p
11. --- 彼女はさらに\VUgras{置時計} \textsuperscript{(57)}、\VUgras{燭台}
\textsuperscript{(58)}、\VUgras{小盆} \textsuperscript{(59)} の埃を払い、それから
\VUgras{鏡} \textsuperscript{(60)} を拭かねばならない。

%%K t06-c1-12-1 p
12. --- 赤子の\VUgras{揺籠} \textsuperscript{(61)}（あの babo (*)）はといえば、
もう用意ができている。

%%K t06-tit-3 sub
\VUsaut{5.0mm}
\VUpk{10.2pt}{40}{客間。}
\VUsaut{3.0mm}

%%K t06-c2-01-1 p
1. --- さて客間である。たいそう美しい\VUgras{部屋}である。右の隅の暖炉には
精巧な\VUgras{小像} \textsuperscript{(1)} と美しい\VUgras{花瓶} \textsuperscript{(3)}
が二つ飾ってあり、天鵞絨の\VUgras{暖炉の帳} \textsuperscript{(2)} (*) が垂れて
いる。

%%K t06-c2-02-1 p
2. --- 炉の火の粉が敷物を焼かぬように、炉の前に銅の\VUgras{火よけ}
\textsuperscript{(4)} が置いてある。

%%K t06-c2-03-1 p
3. --- そのそばに婦人用の洒落た\VUgras{文机} \textsuperscript{(5)} が開いて
いる。\VUgras{家の主婦} \textsuperscript{(23)} はここで手紙を書く。

%%K t06-c2-04-1 p
4. --- 窓には\VUgras{紗の帳} \textsuperscript{(6)} が掛かり、\VUgras{括り紐}
\textsuperscript{(8)} でまとめてある。その紐は\VUgras{掛鉤} \textsuperscript{(7)} に
掛かっている。それに厚い布の帳があり、上に\VUgras{帳の楣} \textsuperscript{(9)} が
ついている。

%%K t06-c2-05-1 p
5. --- 窓の凹みでは\VUgras{花台} \textsuperscript{(11)} が緑の\VUgras{草}
\textsuperscript{(10)} を載せている。見事な棕櫚で、その葉は\VUgras{石油ランプ}
\textsuperscript{(12)} にとどきそうである。

%%K t06-c2-06-1 p
6. --- \VUgras{ランプの笠} \textsuperscript{(13)} はマリが作った。あの愛らしい
\VUgras{娘} \textsuperscript{(14)} で、\VUgras{ピアノ} \textsuperscript{(15)} の前に
いる。\VUgras{ピアノの床几} \textsuperscript{(16)} にまっすぐ坐って曲を稽古して
いる。たいそう腕がよく、たいそう几帳面である。彼女の\VUgras{楽譜の棚}
\textsuperscript{(17)} を見れば分かる。中はきちんと整っている。

%%K t06-tit-4 sub
\VUsaut{5.0mm}
\VUpk{10.2pt}{40}{いたずら者。}
\VUsaut{3.0mm}

%%K t06-c2-07-1 p
7. --- 弟のポールは\VUgras{本} \textsuperscript{(19)} を一冊探している。
\VUgras{本棚} \textsuperscript{(18)} の中をかき回している。この\VUgras{子}
\textsuperscript{(20)} はじっとしていられず、まったく手に負えない。高すぎる所に
ある本にとどこうと本棚にぶら下がって、その上の\VUgras{胸像} \textsuperscript{(21)}
を落としかねない。

%%K t06-c2-08-1 p
8. --- こんな馬鹿をしでかせるのは、母が女友だちと話しこんでいるからである。
その\VUgras{婦人} \textsuperscript{(22)} は数日を共に過ごしに来ている。母は
\VUgras{長椅子} \textsuperscript{(24)} に坐っていて、\VUgras{肘掛椅子}
\textsuperscript{(25)} のそばである。友だちは\VUgras{椅子} \textsuperscript{(27)} に
坐っており、我々にはその\VUgras{背} \textsuperscript{(26)} しか見えない。

%%K t06-c2-09-1 p
9. --- 壁に掛かった大きな\VUgras{額} \textsuperscript{(30)} には一人の親戚の
\VUgras{肖像} \textsuperscript{(29)} が入っている。卓の上に開いた\VUgras{写真帖}
\textsuperscript{(32)} には家のほかの者の\VUgras{写真} \textsuperscript{(33)} が
集めてある。子供たちがさっきまで繰っていたので、\VUgras{椅子}
\textsuperscript{(31)} がまだ卓のまわりに集まっている。

%%K t06-c2-10-1 p
10. --- 磁器の\VUgras{支那の花瓶} \textsuperscript{(34)} は\VUgras{紙切小刀}
\textsuperscript{(35)} のそばにあって、マリが名の日にもらったものである。部屋の
隅を埋める\VUgras{屏風} \textsuperscript{(36)} は叔父が支那から持ち帰ったもので
ある。

%%K t06-c2-11-1 p
11. --- 美しい\VUgras{絵} \textsuperscript{(37)} が\VUgras{壁} \textsuperscript{(42)}
に掛かって部屋を明るくし、天井の\VUgras{吊ランプ} \textsuperscript{(38)} がそれを
照らし、その\VUgras{電球} \textsuperscript{(39)} は夜になるとうれしいほど明るい。
この客間はたいそう心地よく見える。床に敷いたやわらかな\VUgras{敷物}
\textsuperscript{(40)} と、あれほど重宝な\VUgras{電気の呼鈴} \textsuperscript{(41)}
とが、これをすっかり快いものにしている。

%%K t06-tit-5 sub
\VUsaut{3.6mm}
\VUpk{10.2pt}{40}{食堂。}
\VUsaut{3.0mm}

%%K t06-c3-01-1 p
1. --- 食事の鐘が鳴った。マリを除いて一家は卓を囲んでいる。食堂は陶の
\VUgras{暖炉} \textsuperscript{(1)} がよく効いて心地よく暖かい。細い\VUgras{煙突}
\textsuperscript{(2)} がついている。

%%K t06-c3-02-1 p
2. --- この部屋には家具が多くない。壁沿いに、\VUgras{床几} \textsuperscript{(4)}
の上に、美しい小さな\VUgras{手水器} \textsuperscript{(3)} が掛かっている。すぐ
そばが\VUgras{食器棚} \textsuperscript{(5)} で、その上に冷たい料理、菓子、当座
使わぬ器を置く。

%%K t06-c3-03-1 p
3. --- こうして上の段には\VUgras{焼いた鶏} \textsuperscript{(7)}、\VUgras{魚}
\textsuperscript{(8)}、\VUgras{鉢} \textsuperscript{(10)} が見え、鉢には
\VUgras{果物} \textsuperscript{(9)} が盛ってある。下には\VUgras{チーズ} \textsuperscript{(11)}、
\VUgras{パン} \textsuperscript{(12)}、\VUgras{油酢入} \textsuperscript{(13)} が
ある。そこにはサラダを和えるに要るものが揃っている。油、酢、\VUgras{塩}
\textsuperscript{(14)}、\VUgras{胡椒} \textsuperscript{(15)} である。最後に一番下の
段には\VUgras{ケーキ} \textsuperscript{(17)} が一つあり、そのそばが
\VUgras{芥子入} \textsuperscript{(16)} である。

%%K t06-c3-04-1 p
4. --- 食堂の時計は\VUgras{掛時計} \textsuperscript{(20)} といい、たいそう変った
形をしている。木の枠にはめてあり、その枠には\VUgras{晴雨計} \textsuperscript{(19)}
と\VUgras{寒暖計} \textsuperscript{(18)} も入っている。

%%K t06-tit-6 sub
\VUsaut{4.6mm}
\VUpk{10.2pt}{40}{食事の給仕。}
\VUsaut{3.0mm}

%%K t06-c3-05-1 p
5. --- 部屋のまわりの壁には蝋を引いた楢の\VUgras{羽目板} \textsuperscript{(21)}
が張ってある。布の\VUgras{戸の帳} \textsuperscript{(22)} が涼み廊へ通じる戸を
よく塞いでいる。食後に一家はそこへ珈琲を飲みに行く。\VUgras{茶碗}
\textsuperscript{(24)} と\VUgras{受皿} \textsuperscript{(25)} はもう\VUgras{磁器棚}
\textsuperscript{(23)} に並べてある。

%%K t06-c3-06-1 p
6. --- 卓の上、天井には\VUgras{吊ランプ} \textsuperscript{(26)} が取りつけて
ある。その明るい光が夜には食堂じゅうに満ちる。

%%K t06-c3-07-1 p
7. --- では\VUgras{食卓} \textsuperscript{(27)} そのものを見よう。一家が入って
きたとき卓はもう整えてあった。\VUgras{卓掛} \textsuperscript{(28)} の上、めいめい
の席に\VUgras{皿} \textsuperscript{(33)} が一枚、\VUgras{ナプキン}
\textsuperscript{(29)} が一枚、\VUgras{匙} \textsuperscript{(34)} が一本、
\VUgras{肉叉} \textsuperscript{(35)} が一本、\VUgras{小刀} \textsuperscript{(36)} が
一本、\VUgras{杯} \textsuperscript{(32)} が一つ置いてあった。酒の\VUgras{瓶}
\textsuperscript{(30)} と水の\VUgras{硝子瓶} \textsuperscript{(31)} もあった。

%%K t06-c3-08-1 p
8. --- \VUgras{召使} \textsuperscript{(39)} はまず\VUgras{吸物の器}
\textsuperscript{(37)} を運んできた。中の吸物はまだ湯気を立てている。いまは第一の
\VUgras{料理} \textsuperscript{(38)} を給仕している。肉の料理である。それから
すでに名を挙げた料理を運び、最後に\VUgras{菓子}のあと、主人たちが涼み廊へ
行ったのを幸いに、卓を片づけて食堂をもとどおりに整えるであろう。

%%K t06-c3-08-2 p
\textit{注。} --- \textit{食事に関する日常の語はここではあらましだけを挙げた。
この表は「宿屋」（第 13 号）と「市場」（第 15 号）の二つの図表で補われる。}

%%K t06-tit-7 sub
\VUsaut{4.6mm}
\VUpk{10.2pt}{40}{台所。}
\VUsaut{3.0mm}

%%K t06-c4-01-1 p
1. --- 半ば開いた戸から台所を覗くと、絵になるような取り散らかりが見える。
\VUgras{竈} \textsuperscript{(1)} の前では\VUgras{火} \textsuperscript{(3)} がぱちぱち
音を立て、\VUgras{回し焼台} \textsuperscript{(5)} が据えてある。大きな
\VUgras{焼肉} \textsuperscript{(7)} が\VUgras{鉄串} \textsuperscript{(6)} に刺さって
いて、もう焼けかけている。

%%K t06-c4-02-1 p
2. --- いま使ったばかりの\VUgras{鞴} \textsuperscript{(8)} と\VUgras{炭の匙}
\textsuperscript{(9)} が、\VUgras{割木} \textsuperscript{(10)} とともに床に放って
ある。

%%K t06-c4-03-1 p
3. --- 炉の上の棚は整っている。\VUgras{提灯} \textsuperscript{(11)}、
\VUgras{燐寸入} \textsuperscript{(13)} と中の\VUgras{燐寸} \textsuperscript{(12)}、
\VUgras{燭台} \textsuperscript{(14)}、\VUgras{手燭} \textsuperscript{(15)} とその
\VUgras{蝋燭} \textsuperscript{(16)} が、みなそれぞれの場所にある。しかし大きな
\VUgras{炭桶} \textsuperscript{(18)} は、\VUgras{石炭} \textsuperscript{(17)} を
いっぱいに入れて --- その石炭は\VUgras{料理女} \textsuperscript{(23)} がいま穴倉
から取ってきたのである --- 台所の真中に置きっぱなしである。

%%K t06-c4-04-1 p
4. --- じつのところ、この気の毒な女は急がねばならない。昼食を間に合わせる
ためである。いまは\VUgras{煮炊きの竈} \textsuperscript{(19)} の前で\VUgras{汁}
\textsuperscript{(24)} をかき回している。焦がしてはならないのである。

%%K t06-c4-05-1 p
5. --- \VUgras{天火} \textsuperscript{(20)} ではケーキが焼けている。子供たちの
おやつに作ったものである。竈の上には\VUgras{柄杓} \textsuperscript{(21)} と
\VUgras{漉し杓子} \textsuperscript{(22)} が掛かっている。

%%K t06-tit-8 sub
\VUsaut{4.0mm}
\VUpk{10.2pt}{40}{台所の道具。}
\VUsaut{3.0mm}

%%K t06-c4-06-1 p
6. --- 部屋の奥に掛かった\VUgras{鍋の一揃} \textsuperscript{(25)} はたいそう
きれいである。美しい銅の\VUgras{煮鍋} \textsuperscript{(26)}、\VUgras{牛乳壺}
\textsuperscript{(27)}、\VUgras{漉し器} \textsuperscript{(28)}、\VUgras{おろし金}
\textsuperscript{(29)} は念入りに磨いてある。

%%K t06-c4-07-1 p
7. --- \VUgras{塩壺} \textsuperscript{(30)} は小さな戸棚の上、\VUgras{茶瓶}
\textsuperscript{(31)} と\VUgras{油の甕} \textsuperscript{(32)} のそばに置いてある。
\VUgras{抽斗} \textsuperscript{(33)} にはおそらく匙叉と台所の小刀が入っている。

%%K t06-c4-08-1 p
8. --- \VUgras{箒} \textsuperscript{(34)} と\VUgras{羽根の払子} \textsuperscript{(35)}
は隅にある。料理女はいま\VUgras{俎の台} \textsuperscript{(36)} を使ったばかりで、
窓のそばに置きっぱなしにしている。

%%K t06-c4-09-1 p
9. --- 壁には\VUgras{手拭} \textsuperscript{(37)} が一枚掛かっていて、そばが
\VUgras{漏斗} \textsuperscript{(38)} である。台所のこちら側には\VUgras{焼網}
\textsuperscript{(39)}、\VUgras{刻み包丁} \textsuperscript{(40)}、\VUgras{俎}
\textsuperscript{(41)} が置いてある。さらにその上、刻み包丁の上の\VUgras{棚}
\textsuperscript{(42)} には\VUgras{壺} \textsuperscript{(43)}、\VUgras{煮込鍋}
\textsuperscript{(44)} とその\VUgras{蓋} \textsuperscript{(45)}、それに\VUgras{大鍋}
\textsuperscript{(46)} がある。\VUgras{揚鍋} \textsuperscript{(47)} はそのそばに
掛かっている。

%%K t06-c4-10-1 p
10. --- この隅で光っているのは銅の\VUgras{蛇口} \textsuperscript{(48)} だけで
ある。\VUgras{流し} \textsuperscript{(49)} の上には大きな\VUgras{水甕}
\textsuperscript{(50)} が置いてあり、下に小さな戸棚があってたいそう重宝であろう。
そこには\VUgras{酢の樽} \textsuperscript{(51)} とその\VUgras{小さな蛇口}
\textsuperscript{(52)}、\VUgras{塵箱} \textsuperscript{(53)}、\VUgras{汚れた器}
\textsuperscript{(54)} が入れてある。

%%K t06-tit-9 sub
\VUsaut{4.0mm}
\VUpk{10.2pt}{40}{台所に置いてあるほかのもの。}
\VUsaut{3.0mm}

%%K t06-c4-11-1 p
11. --- あの大きな\VUgras{木の盥} \textsuperscript{(55)} は\VUgras{野菜}
\textsuperscript{(58)} を洗うためのもの、それに鋳鉄の\VUgras{釜}
\textsuperscript{(56)} は\VUgras{敷瓦の床} \textsuperscript{(57)} に据えてあるが、
これもおそらくあの戸棚のものであろう。

%%K t06-c4-12-1 p
12. --- 料理女に竈が足りぬことはない。こちらの卓の隅に\VUgras{瓦斯の竈}
\textsuperscript{(59)} があり、その上の小さな鍋で湯が沸いている。立ちのぼる
\VUgras{湯気} \textsuperscript{(60)} が、もう煮立っているに違いないと教えている。
この湯はおそらく珈琲のためであろう。卓の反対の端に\VUgras{珈琲沸し}
\textsuperscript{(64)} と\VUgras{珈琲挽} \textsuperscript{(63)} があるからである。

%%K t06-c4-13-1 p
13. --- 瓦斯の竈は\VUgras{瓦斯の口} \textsuperscript{(62)} につないである。長い
\VUgras{ゴムの管} \textsuperscript{(61)} を使ってである。

%%K t06-c4-14-1 p
14. --- 料理女を手伝いに来る\VUgras{通いの女} \textsuperscript{(66)} が夕食の
野菜を用意している。皮を剥き洗い終えたら、煮る時が来るまでそれを\VUgras{器}
\textsuperscript{(65)} に入れておくのである。

%%K t06-tit-10 sub
\VUsaut{4.0mm}
{\centering\fontsize{10.2pt}{10.2pt}\selectfont\textls[40]{\scshape 湯殿。} \textit{（行水の部屋。）}\par}
\VUsaut{3.0mm}

%%K t06-c5-01-1 p
1. --- この家の主な部屋を残らず知るには、あと一度だけ覗けばよい。湯殿の
設えを見るのである。長くはかかるまい。広くはないからで、すぐに分かるとおり
\VUgras{湯槽} \textsuperscript{(1)} にはよい\VUgras{湯沸器} \textsuperscript{(2)} が
ついており、\VUgras{石鹸皿} \textsuperscript{(3)} は入る者の手の届くところにあり、
\VUgras{手拭掛} \textsuperscript{(5)} があれば\VUgras{手拭} \textsuperscript{(4)} は
たやすく乾くであろう。

%%K t06-c5-02-1 p
2. --- この部屋は壁に\VUgras{釉薬の敷瓦} \textsuperscript{(6)} が貼ってある。
だから\VUgras{シャワー} \textsuperscript{(7)} を取りつけても、使うときに跳ねる水が
壁を傷めることを恐れずにすんだ。足を洗う小さな亜鉛の\VUgras{盥}
\textsuperscript{(8)} が備えを整えている。
\VUsaut{6.0mm}
\VUfilet{22mm}
